\section{Results}

\subsection{The null experiment isolated deterministic component coupling from contingent total effects}

All \PermutationRows{} permutation-level cells were present, and the numerical and integrity gates passed. At the maximum regression budget, the squared mean-gap contribution favored response-aware selection in 100\% of endpoint permutations for ESOL and FreeSolv under both scaffold conventions. Under the shared-acyclic convention, the mean squared-gap contributions to MSE were \ESGapSqMean{} for ESOL and \FSGapSqMean{} for FreeSolv. The corresponding mean test-variance contributions were \ESVarianceMean{} and \FSVarianceMean{}. Consequently, total RMSE effects were much less stable: \ESNullRMSEMean{} with central 95\% null interval [\ESNullRMSELow{}, \ESNullRMSEHigh{}] for ESOL and \FSNullRMSEMean{} [\FSNullRMSELow{}, \FSNullRMSEHigh{}] for FreeSolv. Both total-effect intervals spanned zero.

The budget trajectories in Figures~\ref{fig:null-rmse} and \ref{fig:null-anatomy} show why the distinction matters. The objective-aligned squared-gap component was consistently positive, but collateral test-variance changes sometimes reinforced and sometimes offset it. A larger candidate pool increased the opportunity to find endpoint-aligned partitions, yet mean total RMSE effects did not follow a universal monotone trajectory. The null experiment therefore supports a component-level mechanism rather than a blanket claim that response-aware selection must reduce headline RMSE.

For classification at 300 draws, mean Brier effects ranged from \ClassBrierMin{} to \ClassBrierMax{} and mean log-loss effects from \ClassLogLossMin{} to \ClassLogLossMax{} across datasets. Constant-score ROC--AUC effects were exactly zero wherever defined. The squared prevalence-gap component was nonnegative by construction, but changes in Bernoulli test variance, test prevalence, and entropy produced small dataset-dependent total effects. ClinTox, for example, had negative mean Brier and log-loss effects despite a reduced prevalence gap. Full classification budget curves and decomposition tables are in Supporting Information.

\begin{figure}[!htbp]
\centering
\includegraphics[width=0.98\textwidth]{figures/Figure_1.pdf}
\caption{Endpoint-permutation null RMSE effects across nested requested-draw budgets for ESOL and FreeSolv under shared-acyclic and singleton conventions. Lines show null means and shaded bands show central 2.5th--97.5th percentiles across 200 permutation-level statistics.}
\label{fig:null-rmse}
\end{figure}

\begin{figure}[!htbp]
\centering
\includegraphics[width=0.98\textwidth]{figures/Figure_2.pdf}
\caption{Decomposition of the null MSE effect into the squared train--test mean-gap component and collateral test-variance component across nested search budgets.}
\label{fig:null-anatomy}
\end{figure}

\begin{figure}[!htbp]
\centering
\includegraphics[width=0.98\textwidth]{figures/Figure_3.pdf}
\caption{Classification null effects across nested requested-draw budgets. Constant-score ROC--AUC is invariant when both classes are represented, whereas Brier score, log loss, and average precision respond to prevalence and composition through different mechanisms.}
\label{fig:null-classification}
\end{figure}

\subsection{Real endpoint--scaffold organization amplified the mean-only effect in three of four matched analyses}

The exploratory bridge produced a strong separation between empirical and permuted endpoint effects for the primary scaffold convention (Table~\ref{tab:bridge}). ESOL's empirical mean-only RMSE effect was \ESPrimaryObserved{}, compared with null mean \ESPrimaryNullMean{} and central null interval [\ESPrimaryNullLow{}, \ESPrimaryNullHigh{}]. FreeSolv's empirical effect was \FSPrimaryObserved{}, compared with \FSPrimaryNullMean{} [\FSPrimaryNullLow{}, \FSPrimaryNullHigh{}]. Both empirical values exceeded all 200 randomized values, giving the smallest attainable plus-one upper-tail probability, \BridgeMinP{}.

At the matched 5,000-draw singleton budget, ESOL's empirical effect, \ESSingletonObserved{}, also exceeded all randomized values. In contrast, singleton FreeSolv's empirical effect, \FSSingletonObserved{}, lay within its null distribution (null mean \FSSingletonNullMean{}; percentile \FSSingletonPercentile{}\%). This pattern mirrors the empirical scaffold-semantic sensitivity: real-data amplification remained for ESOL but disappeared for FreeSolv when acyclic molecules no longer formed one shared scaffold group.

Because the bridge was added after null completion, it is evidence-generating rather than confirmatory. Nevertheless, it sharpens the mechanism: endpoint-aware selection always reduces an aligned squared-gap component, while strong observed headline effects require the actual response to be organized across the available scaffold candidate space in a way that exceeds random reassignment.

\begin{table}[!htbp]
\centering
\caption{Exploratory empirical--null bridge for the mean-only RMSE effect. Null intervals are central 2.5th--97.5th percentiles across 200 endpoint permutations after averaging 20 partition seeds within each permutation.}
\label{tab:bridge}
\input{generated/bridge_table.tex}
\end{table}

\subsection{Empirical learned-model effects remained protocol-dependent}

The learned-model results independently confirmed that benchmark choices can alter apparent performance. None of 12 classification dataset--model cells met the corrected response-aware advantage criterion. All six shared-acyclic regression cells did so, but the regression pattern attenuated strongly under singleton acyclic semantics. The empirical mean-only effect exceeded the corresponding learned-model mean effect in all six primary regression cells.

The statistical interpretation remains conservative. No corrected classification advantage does not imply equivalence, and the primary regression advantages do not imply that response-aware selection is a preferred deployment split. They show that a searched benchmark can change the difficulty and composition of the held-out population.

\begin{figure}[!htbp]
\centering
\includegraphics[width=0.96\textwidth]{figures/Figure_4.pdf}
\caption{Empirical learned-model paired effects. Panel A shows the 12 classification AUC effects. Panel B shows the six regression RMSE effects. Points are means over 20 unique partition pairs; bars are paired-bootstrap 95\% intervals. Positive values favor response-aware selection. Corrected decisions were 0/12 classification and 6/6 shared-acyclic regression advantages.}
\label{fig:empirical-primary}
\end{figure}

\begin{figure}[!htbp]
\centering
\includegraphics[width=0.96\textwidth]{figures/Figure_5.pdf}
\caption{Acyclic-scaffold semantic sensitivity. The shared-acyclic primary and singleton sensitivity kept endpoints, models, partition seeds, and exact-size pairing logic but imposed different structural meanings for acyclic molecules. Shared-acyclic effects attenuated for ESOL and reversed for FreeSolv under singleton semantics.}
\label{fig:acyclic-primary}
\end{figure}
